La proportionnalité est la première structure du programme : deux grandeurs sont
proportionnelles quand l'une est un multiple fixe de l'autre, et tout ce qui suit — quatrième
proportionnelle, pourcentages, échelles, vitesse moyenne — en découle par calcul.

Le chapitre introduit ensuite la notion de fonction, et c'est là que la formalisation apporte
quelque chose d'inattendu. Une fonction, en Lean comme en mathématiques, est un objet total :
à chaque nombre correspond une image. Or l'usage scolaire parle volontiers de \og{} la
fonction $1/x$ \fg{} sans domaine, et lit des antécédents sur un graphique. Les énoncés
retenus ici sont ceux qui portent sur l'objet et non sur sa représentation : une lecture
graphique n'est pas une proposition.
